\section{直和,自由模,分次模的张量代数}

记号约定:$A$是交换环,$M$是$A$-模.

\begin{proposition}{模的直和的张量代数}{}
	设$\left(M_{\lambda}\right)_{\lambda\in\Lambda}$是一族$A$-模,$M=\bigboxplus\limits_{i\in I}M_{i}$.对每个$n\in\N$和$\lambda_{1},\ldots,\lambda_{n}\in\Lambda$,记
	\begin{align*}
		T_{\lambda_{1}\ldots\lambda_{n}}:M_{\lambda_{1}}\otimes\ldots\otimes M_{\lambda_{n}}\to\mathsf{T}^{n}\left(M\right)
	\end{align*}
	为典范单射,则$\mathsf{T}\left(M\right)=\bigboxplus\limits_{n\in\N}\bigboxplus\limits_{\lambda_{1},\ldots,\lambda_{n}\in\Lambda}\im\left(T_{\lambda_{1}\ldots\lambda_{n}}\right)$.
\end{proposition}

\begin{theorem}{自由模的交换代数}{}
	设$A$是交换环,$M$是自由$A$-模,$\left(e_{\lambda}\right)_{\lambda\in\Lambda}$是$M$的基.对$\Lambda$的每个有限序列$\mathbf{s}=\left(\lambda_{i}\right)_{i=1}^{n}$,定义
	\begin{align*}
		e_{\mathbf{s}}=e_{\lambda_{1}}\otimes\ldots\otimes e_{\lambda_{n}}.
	\end{align*}
	约定:$\mathbf{s}$是空序列时$e_{\mathbf{s}}=1$,那么$\left(e_{\mathbf{s}}\right)_{\mathbf{s}\in\Lambda^{n}}$是$\mathsf{T}\left(M\right)$的基.
\end{theorem}

\begin{remark}{}{}
	$\left(e_{\mathbf{s}}\right)_{\mathbf{s}\in\Lambda^{n}}$是一族齐次元,它在$\mathsf{T}\left(M\right)$中的乘法表满足$\forall\mathbf{s},\mathbf{t}\in\Lambda^{n}\left(e_{\mathbf{s}}e_{\mathbf{t}}=e_{\mathbf{s}\mathbf{t}}\right)$,其中$\mathbf{s}\mathbf{t}$是$\mathbf{s}$和$\mathbf{t}$并接得到的新序列.因此,配备了这个基和上述乘法法则的$\mathsf{T}\left(M\right)$典范同构于$\Libas_{A}\left(\Lambda\right)$.特别地,如果$E$是$A$-代数,那么对每个$f\in\Hom_{\mathsf{Set}}\left(\Lambda,E\right)$,存在唯一的$\bar{f}\in\Hom_{A-\mathsf{Alg}}\left(\mathsf{T}\left(M\right),E\right)$满足$\forall\lambda\in\Lambda\left(\bar{f}\left(e_{\lambda}\right)=f\left(\lambda\right)\right)$.
\end{remark}

\begin{corollary}{投射模的张量代数也是投射模}{}
	如果$M$是投射模,那么$\mathsf{T}\left(M\right)$也是投射模.
\end{corollary}

\begin{proposition}{分次模的张量代数}{}
	设$\Delta$是交换幺半群(加法记号),$M$是带有分次$\left(M_{\alpha}\right)_{\alpha\in\Delta}$的$A$-模.对每个$\left(\alpha,n\right)\in\Delta\times\N$,定义
	\begin{align*}
		\mathsf{T}^{\alpha,n}\left(M\right)=\bigboxplus\limits_{\substack{\alpha_{1},\ldots,\alpha_{n}\in\N \\ \alpha_{1}+\ldots+\alpha_{n}=\alpha}}M_{\alpha_{1}}\otimes\ldots\otimes M_{\alpha_{n}}.
	\end{align*}
	\begin{enumerate}
		\item $\left(\mathsf{T}^{\alpha,n}\left(M\right)\right)_{\alpha\in\Delta,n\in\N}$是$\mathsf{T}\left(M\right)$上的$\Delta\times\N$型分次.
		\item $\mathsf{T}\left(M\right)$上既与$\mathsf{T}\left(M\right)$的$A$-代数结构相容,又诱导$\mathsf{T}^{1}\left(M\right)$上的$\Delta$型分次的分次只有$\left(\mathsf{T}^{\alpha,n}\left(M\right)\right)_{\alpha\in\Delta,n\in\N}$.
	\end{enumerate}
\end{proposition}

























